\documentclass[10pt,a4paper]{amsart}
\usepackage[margin=19mm]{geometry}
\usepackage{amsmath,amssymb,mathtools}
\usepackage[scr=rsfs,cal=euler]{mathalfa}
\usepackage{xfrac}
\usepackage[unicode,hidelinks,bookmarks=false]{hyperref}
\newcommand{\ie}{i.e.\ }
\newcommand{\cf}{cf.\ }
\newcommand{\resp}{respectively\ }
\newcommand{\Subsection}{\subsection}
\providecommand{\nobreakdash}{\nobreak\hbox{-}}
\setlength{\emergencystretch}{2em}
\multlinegap0pt
\usepackage{bm}
\usepackage{bbm}
\usepackage{dsfont}
\usepackage{xspace}
\usepackage{cancel}
\usepackage{mathtools}
\usepackage{amsthm}
\makeatletter
\@ifundefined{theorem}{\newtheorem{theorem}{Theorem}}{}
\@ifundefined{proposition}{\newtheorem{proposition}[theorem]{Proposition}}{}
\makeatother
\makeatletter
\expandafter\ifx\csname proof\endcsname\relax
\newenvironment{proof}[1][Proof]{\noindent\textbf{#1.} }{\ \rule{0.5em}{0.5em}}
\fi
\makeatother
\title[Constructive Notes on Locally Convex Spaces]{Constructive Notes on Locally Convex Spaces}
\author{Douglas S. Bridges}
\date{Source: arXiv:2606.12746v1 (CC BY 4.0)}
\begin{document}
\maketitle

\noindent\emph{Source and licence.} Constructive Notes on Locally Convex Spaces. Version arXiv:2606.12746v1 (CC BY 4.0), distributed under the Creative Commons Attribution 4.0 International licence, as recorded in the arXiv licence metadata for this exact version and shown on the abstract page https://arxiv.org/abs/2606.12746v1. Full rights notice: licenses/arxiv-2606.12746-CC-BY-4.0.txt.\par\smallskip

\noindent\textit{Editorial scope.} Counted unit: the Proposition whose source label is \texttt{p11} in the article (source0.tex lines 639--645, source label \texttt{p11}), delivered together with the article's own complete original proof of it (source0.tex lines 647--734, one \texttt{proof} environment of 88 lines). No mathematics is written, repaired or re-derived: statement and proof are the article's own text line for line. Required context retained uncounted: none. Every definition, display and label of those lines is retained unchanged. Omitted from this package and not redistributed: 1--638, 646--646, 735--1103. Nothing inside a retained excerpt is omitted or altered: every retained body is delivered exactly as the article's lines are written, so no omission receipt of any kind accompanies this package. Citations: the retained excerpts carry no citation command at all, so no bibliography entry of the article is transcribed and no reference is redistributed. Reference adaptation: none was needed and none was made; every reference inside the delivered unit resolves inside it, so no unresolved reference or citation is produced. Printed slips kept literal: no printed word, symbol, hypothesis or number of the counted statement or of its proof was altered. Layout and class: the article's own class is \texttt{article} with packages including amsmath, amssymb, euler, amsfonts, graphicx; the wrapper is the lane's pinned \texttt{amsart} editorial wrapper at 10pt on A4, which restates the theorem-like environments the retained text needs and adds the article's own macro definitions that the retained text needs (none), with extra wrapper packages (bm, bbm, dsfont, xspace, cancel, mathtools). The delivered numbering is the unit's own. Assets: no retained span contains a figure environment, a graphics-inclusion command, a PDF-page inclusion, a table or a TikZ picture; no image file is copied into this package, the package declares no asset dependency and no image file is redistributed with it. Licence: version arXiv:2606.12746v1 of this article is distributed under the Creative Commons Attribution 4.0 International licence, as recorded in the arXiv licence metadata for this exact version and shown on the abstract page https://arxiv.org/abs/2606.12746v1. A full rights notice is delivered at licenses/arxiv-2606.12746-CC-BY-4.0.txt. Characters: every retained excerpt is pure ASCII, so the delivered engine, XeLaTeX with its default Latin Modern text font, reports no missing character.
\begin{proposition}
\label{p11}Let $S$ be a totally bounded subset of the locally convex space
$\left(  X,\left(  p_{i}\right)  _{i\in I}\right)  $, and $x_{0}\in S$. Let
$F$ be a finitely enumerable subset of $I$, and $r>0$. Then there exists a
closed $F$-totally bounded subset $K$ of $S$ such that $S\cap B^{F}%
(x_{0},r)\subset K\subset S\cap\overline{B}^{F}(x_{0},8r)$.
\end{proposition}

\begin{proof}
With $S_{1}\equiv\{x_{0}\},$ we construct inductively a sequence
$(S_{n})_{n\geqslant1}$ of finitely enumerable subsets of $S$ such that

\begin{itemize}
\item[(a)] $\rho^{F}(x,S_{n})<2^{-n+1}r$ for each $x$ in $S\cap B^{F}%
(x_{0},r)$, and

\item[(b)] \smallskip$\rho^{F}(x,S_{n})<2^{-n+3}r$ for each $x$ in $S_{n+1}.$
\end{itemize}

%

%TCIMACRO{\TeXButton{noindent}{\noindent}}%
%BeginExpansion
\noindent
%EndExpansion
To that end, assume that $S_{1},\ldots,S_{n}$ have been constructed with the
appropriate properties, and let $\{x_{1},\ldots,x_{N}\}$ be an $(F,2^{-n}%
r)$-approximation to $S.$ Write $\{1,\ldots,N\}$ as a union of subsets $A$ and
$B$ such that%
\begin{align*}
\rho^{F}(x_{k},S_{n})  &  <2^{-n+3}r\quad\mathrm{if}\,\,k\in A,\\
\rho^{F}(x_{k},S_{n})  &  >2^{-n+2}r\quad\mathrm{if}\,\,k\in B.
\end{align*}
Let%
\[
S_{n+1}\equiv\{x_{k}:k\in A\}.
\]
Clearly, $S_{n+1}$ satisfies (b). Let $x$ be any point of $S\cap B^{F}%
(x_{0},r).$ By the induction hypothesis, there exists $y$ in $S_{n}$ with
$\sum_{i\in F}p_{i}(x-y)<2^{-n+1}r.$ Choosing $k$ in $\{1,\ldots,N\}$ such
that $\sum_{i\in F}p_{i}(x-x_{k})<2^{-n}r,$ we have
\begin{align*}
\rho^{F}(x_{k},S_{n})  &  \leqslant\sum_{i\in F}p_{i}(x_{k}-y)\\
&  \leqslant\sum_{i\in F}p_{i}(x-x_{k})+\sum_{i\in F}p_{i}(x-y)\\
&  <2^{-n}r+2^{-n+1}r\\
&  <2^{-n+2}r.
\end{align*}
Thus $k\notin B$, so $k\in A$, which is therefore (inhabited and) finitely
enumerable. Moreover,
\[
\rho^{F}(x,S_{n+1})\leq\sum_{i\in F}p_{i}(x-x_{k})<2^{-(n+1)+1}r,
\]
so $S_{n+1}$ satisfies the appropriate instance of (a). This completes the
inductive construction. Letting $K$ be the closure of $\bigcup\limits_{n=1}%
^{\infty}S_{n}$ in $X$, we see from (a) that $S\cap B^{F}(x_{0},r)\subset K.$

Next we prove:

\begin{itemize}
\item[(*)] \emph{If }$m,n$\emph{\ are positive integers with }$m\geq n$\emph{,
and }$y\in S_{m}$\emph{, then there exists }$y_{n}\in S_{n}$\emph{\ such that
}$\sum_{i\in F}p_{i}(y-y_{n})<2^{-n+4}r$\emph{.}
\end{itemize}

%

%TCIMACRO{\TeXButton{noindent}{\noindent}}%
%BeginExpansion
\noindent
%EndExpansion
Given $y\in S_{m}$ and letting $y_{m}=y$, we see from (b) that for each $k$
with $n\leq k\leq m-1$ there exist points $y_{k}\in S_{k}$ such that
$\sum_{i\in F}p_{i}(y_{k+1}-y_{k})<2^{-k+3}r$. Then%
\[
\sum_{i\in F}p_{i}(y-y_{n})\leqslant\sum_{k=n}^{m-1}\sum_{i\in F}p_{i}%
(y_{k+1}-y_{k})<\sum_{k=n}^{\infty}2^{-k+3}r=2^{-n+4}r.
\]
In particular taking $n=1$, we see that $\rho^{F}(y,\left\{  x_{0}\right\}
)=\sum_{i\in F}p_{i}(y-y_{1})<2^{3}r~$for each $y\in\bigcup\limits_{i=1}%
^{\infty}S_{i}$; whence $\bigcup\limits_{n=1}^{\infty}S_{n}\subset S\cap
B^{F}(x_{0},8r)$ and therefore $K\subset\overline{B}^{F}(x_{0},8r)$. On the
other hand, given $x\in K$ and a positive integer $n$, we can find $m$ and
$y\in S_{m}$ such that $\sum_{i\in F}p_{i}(x-y)<2^{-n+4}r$. If $m<n$, then
$y\in\bigcup\limits_{k=1}^{n}S_{k\text{.}}$. If $m\geq n$, then by the remark
preceding this proposition,%
\[
\rho^{F}(x,S_{n})\leqslant\sum_{i\in F}p_{i}(x-y)+\rho^{F}(y,S_{n}%
)<2^{-n+4}r+2^{-n+4}r=2^{-n+5}r,
\]
the second inequality using (*). Thus there exists $z\in S_{n}\subset
\bigcup\limits_{k=1}^{n}S_{k\text{. }}$ such that $\sum_{i\in F}%
p_{i}(x-z)<2^{-n+5}r$. Putting together the two alternatives for $m$, we now
see that the finitely enumerable set $\bigcup\limits_{k=1}^{n}S_{k\text{. }}%
$is a $\left(  F,2^{-n+5}r\right)  $-approximation to $K$. Since $n$ is
arbitrary, we conclude that $K$ is $F$-totally bounded.
\end{proof}

\end{document}
